% !TeX program = xelatex
% !TeX root = draft_proposal.tex
% Calibri requires XeLaTeX or LuaLaTeX + system font (e.g. Microsoft Office / Windows fonts on macOS).
\documentclass[12pt,a4paper]{article}

% Font note:
% - This document is meant to compile reliably on systems without MS Office fonts.
% - XeLaTeX handles Unicode (e.g., Ö, Ş) without extra inputenc settings.

\usepackage[margin=1in]{geometry}
\usepackage{setspace}
\usepackage{enumitem}
\usepackage{amsmath}
\usepackage{cite}
\usepackage{booktabs}
\usepackage{url}
\usepackage{graphicx}
\usepackage[hidelinks]{hyperref}

\setstretch{1.15}

\begin{document}

% --- Cover page (centered; logo at bottom) ---
\thispagestyle{empty}
\begin{center}
    \vspace*{1.25em}
    {\LARGE\bfseries ENS 491 -- Graduation Project (Design)}\\[0.85em]
    {\large\bfseries Proposal}

    \vfill

    {\bfseries Project Title}\\[0.55em]
    {\normalsize AI-Driven Uncertainty-Aware Visualization for Decision Making Under Stress}

    \vspace{2.25em}

    {\bfseries Group Members}\\[0.55em]
    {\normalsize Mustafa Bozyel (\#32417), Ömer Mert Özel (\#31998), Senih Kırmaç (\#32426)}

    \vspace{1.35em}

    {\bfseries Supervisor(s)}\\[0.55em]
    {\normalsize Selim Balcısoy}

    \vspace{1.35em}

    {\bfseries Date}\\[0.55em]
    {\normalsize April 17, 2026}

    \vfill

    \includegraphics[width=0.58\textwidth]{images/sabanci-logo.jpeg}
\end{center}
\clearpage

\section{ABSTRACT}

Many decision-support systems present uncertain, incomplete, or ambiguous data as if it were precise. Under stress---for example in emergency response, security monitoring, or other time-critical operations---this can encourage overconfident or unsafe decisions. The overarching graduation project targets an AI-driven approach that explicitly models uncertainty and communicates it through interactive visualizations, while remaining cognitively manageable for operators and supporting interpretability, latency, and trust constraints.

Within that framing, this proposal concentrates on a concrete, research-ready slice: \textbf{visualizing uncertainty in GPS and location-based data}. Location fixes are often shown as a single point or a uniform ``blue dot,'' even though sensor noise, multipath, map-matching error, and semantic ambiguity mean that true position may lie in a region rather than at one coordinate. Our goal is not primarily to eliminate uncertainty but to \textbf{represent it honestly} through map-based encodings that users can interpret. We survey prior work on degraded-GPS navigation displays, blue-dot-style glyphs, and GIS-oriented GPS uncertainty visualization \cite{ranasinghe2019visualising,burigat2011pedestrian,mckenzie2016assessing,hegarty2016where,lodha2002visualization}, and we outline design directions (e.g., probability-weighted boundaries, internal density) plus an implemented web prototype that loads Android GNSS Logger exports and compares point, heatmap, and accuracy-ring views. We also relate our trajectory-focused use case to broader issues in dense mobility visualization and GPS preprocessing \cite{Zheng2015,Sobral2019,Fadloun2022,Feng2021,MillardBall2019}. The expected outcome is a clearer design space for uncertainty-aware geospatial interfaces that can later be integrated into stress-aware decision-support scenarios.

\section{INTRODUCTION}

\subsection*{Background and problem}

The official project theme is \textit{AI-Driven Uncertainty-Aware Visualization for Decision Making Under Stress}. Software frequently masks uncertainty, which is especially risky when operators must act quickly. Stakeholders include end users, system designers, and safety authorities; constraints include interpretability, latency, and trust.

We scope the \textbf{first major technical track} to \textbf{GPS and location uncertainty visualization}: how to show that a reported position is a distribution or region, not a fact at a pixel. Example: a transaction may geolocate to a mall while the true place is one of many indoor venues; a GPS fix may be drawn as one coordinate while the device could lie anywhere inside an error ellipse. Our guiding question is: how can uncertainty in GPS and location data be visualized in a way that is honest, interpretable, and useful for decision-making?

Location uncertainty arises from sensor noise, signal loss, multipath in urban canyons, latency, map-matching and aggregation errors, and model-based inference. We distinguish positional, semantic, model, and temporal uncertainty categories to structure the design space.

\subsection*{Literature survey and limitations of prior work}

\textbf{Navigation under degraded GPS.} Ranasinghe et al.\ evaluated novel uncertainty visualizations for pedestrian navigation in the field; landmark-oriented designs helped users judge true location under poor GPS \cite{ranasinghe2019visualising}. Burigat and Chittaro compared circular uncertainty regions and street-based encodings; coloring streets reduced workload and was rated helpful \cite{burigat2011pedestrian}. These studies show that \textbf{encoding choice changes performance} but often focus on specific navigation tasks rather than general decision dashboards.

\textbf{Blue-dot and judgment tasks.} McKenzie et al.\ compared uniform circles, Gaussian fades, and centroid dots for positional uncertainty judgments; uniform circles yielded faster, distribution-consistent judgments in their setting \cite{mckenzie2016assessing}. Hegarty et al.\ showed that users combine uncertain cues in task-dependent ways \cite{hegarty2016where}. \textbf{Shortcoming:} uniform disks do not communicate \emph{where} probability mass is higher inside the disk---a gap our proposed boundary- and interior-density ideas target.

\textbf{GIS and trajectories.} Lodha et al.\ modeled horizontal, speed, and direction error for moving GPS objects and visualized diffused representations over imagery \cite{lodha2002visualization}. Broader trajectory and mobility literature emphasizes preprocessing (filtering, segmentation, map matching), clutter, and density visualization \cite{Zheng2015,Sobral2019,Fadloun2022,Feng2021,MillardBall2019}. \textbf{Shortcoming:} classical density maps can hide individual structure; planar KDE may misrepresent network-constrained motion.

\textbf{Commercial products.} Consumer maps use proprietary uncertainty metaphors; details are not open for replication. Our work emphasizes \textbf{transparent, research-grade} encodings we can evaluate and publish.

\subsection*{Patents}

We have not identified a specific patent that blocks our \textbf{research prototype} on uncertainty visualization for logged GNSS data. Commercial map and device vendors hold broad portfolios; we avoid copying protected assets and rely on published methods and our own implementations. A formal freedom-to-operate review can be added if the work moves toward a product.

\subsection*{Motivation}

Making uncertainty visible supports more realistic trust calibration, aligns with human-centered AI goals, and provides a stepping stone toward stress-aware dashboards where misleading precision would be especially harmful. The GPS slice is tractable, data is easy to collect (e.g., GNSS Logger), and findings transfer to other uncertain spatial evidence.

\section{PROPOSED SOLUTION AND METHODS}

\subsection*{What we propose}

We propose an iterative pipeline: (1) expand literature on uncertainty visualization and dense mobility; (2) specify visual encodings for GPS uncertainty (including concepts beyond the uniform blue dot); (3) implement and refine interactive prototypes (current work: a web tool with point, grid heatmap, kernel-density heatmap, and accuracy-ring modes on real \texttt{Fix} logs); (4) evaluate interpretability, cognitive load, and honesty of communication, with an eye toward later integration into broader stress-aware, AI-assisted decision support.

\subsection*{Scope}

\textbf{In scope:} visualization of positional uncertainty from smartphone/GNSS logs; exploratory comparison of encodings; preprocessing basics (sorting, cleaning, derived fields such as accuracy columns); relation to trajectory clutter and density visualization literature.

\textbf{Out of scope (initially):} full probabilistic localization stack for production navigation; real-time emergency dispatch integration; rigorous user studies at scale (planned later); complete modeling of stress physiology---we treat stress as motivation for the \emph{importance} of honest uncertainty, while the thesis work grounds in GPS visualization.

\subsection*{Why this is a ``complex problem''}

\begin{itemize}[leftmargin=1.5em]
    \item \textbf{Multiple constraints:} perceptual limits, map readability, mobile latency, trust, and correctness of uncertainty communication simultaneously.
    \item \textbf{No single obvious solution:} optimal glyphs depend on task (navigation vs.\ monitoring vs.\ forensics) \cite{hegarty2016where}.
    \item \textbf{Research-based:} draws on HCI, visualization, GIS, and trajectory mining \cite{Zheng2015,Sobral2019,Fadloun2022}.
    \item \textbf{Non-routine engineering:} combining map interaction, density estimation, and reported accuracy fields in one tool is integration-heavy.
    \item \textbf{Consequences:} misleading location confidence affects safety and privacy (see Ethical Issues).
    \item \textbf{System-level:} data ingestion, UI, evaluation, and future AI layers must align.
\end{itemize}

\subsection{Objectives/Tasks}

\begin{center}
\small
\begin{tabular}{@{}p{0.42\textwidth}p{0.48\textwidth}@{}}
\toprule
\textbf{Objective / task} & \textbf{Intended result} \\
\midrule
Literature and design space & Synthesized requirements and encoding catalog \\
GNSS data collection \& logs & Reproducible \texttt{Fix}-only datasets for tests \\
Web prototype (Leaflet/D3) & Runnable comparison of visualization modes \\
Concept designs (boundary/density) & Illustrations and design rationale \\
Evaluation plan & Metrics (interpretability, load, trust); pilot study design \\
Report and advisor alignment & ENS 491/492 deliverables updated each term \\
\bottomrule
\end{tabular}
\end{center}


\subsection{Realistic Constraints}

\begin{itemize}[leftmargin=1.5em]
    \item \textbf{Environmental:} Field logging has negligible environmental impact; we avoid unnecessary travel for data collection.
    \item \textbf{Social/Privacy:} Location traces are sensitive; we use consented, locally stored logs and avoid publishing raw coordinates at identifiable resolution in public repos (see Ethics).
    \item \textbf{Health and Safety:} Outdoor walking data collection should follow normal pedestrian safety; no hazardous materials.
    \item \textbf{Manufacturability / deployment:} Prototype is a browser application; depends on modern browsers and JavaScript; GitHub Pages hosting is adequate for demos.
\end{itemize}

\subsection{Engineering/Scientific Standards}

We align with good practice for \textbf{empirical visualization and HCI research} (transparent stimuli, documented parameters) and with \textbf{geospatial data quality} principles (metadata such as \texttt{AccuracyMeters}, honest labeling of synthetic vs.\ real data). Where map matching or road networks become central, established GIS and transportation methods apply \cite{MillardBall2019}. We do not claim compliance with a single mandatory industry code; the advisor may require additional standards (e.g., accessibility) as the project evolves.

\subsection{Sustainability and Life Cycle Assessment}

\noindent\textbf{Sustainability:} Open formats (text logs) aid long-term reproducibility; documentation reduces dependency on single maintainers.

\begin{itemize}[leftmargin=1.5em]
    \item \textbf{Functional unit:} One complete analysis session consisting of (i) importing one GNSS Logger \texttt{Fix}-only trace (e.g., 5--30 minutes of walking/driving), (ii) rendering map-based uncertainty views (points/heatmap/accuracy rings), and (iii) generating the deviation summary versus an OSM reference polyline.
    \item \textbf{System boundary (gate-to-gate, software-focused):} Smartphone GNSS logging (existing devices), local preprocessing in browser/JavaScript, and optional static web hosting for demos (e.g., GitHub Pages). We exclude device manufacturing and user travel beyond normal commuting/walking because they are not driven by the software itself.
    \item \textbf{Energy consumption estimate:} A typical session on a laptop is expected to be CPU-light and short. Assuming a 30 W average laptop draw for 10 minutes of active analysis, energy use is
    \[
      E \approx 30~\text{W} \times \frac{10}{60}~\text{h} \approx 5~\text{Wh} = 0.005~\text{kWh}.
    \]
    \item \textbf{Carbon footprint (kg CO$_2$ eq):} Using a conservative grid intensity of 0.4 kg CO$_2$/kWh, the operational footprint per session is:
    \begin{align*}
      \text{CO}_2 &\approx 0.005~\text{kWh} \times 0.4~\frac{\text{kg CO}_2}{\text{kWh}} \\
                 &\approx 0.002~\text{kg CO}_2 \;\; (\sim 2~\text{g}).
    \end{align*}
    Static hosting for demos is negligible at this scale; we primarily rely on local execution.
    \item \textbf{Comparison to a conventional alternative:} A ``blue-dot only'' or ``single-point'' map has similar compute cost but can increase downstream costs (time, risk, rework) by encouraging overconfident interpretation. Our approach aims to reduce decision error without increasing operational energy meaningfully.
\end{itemize}

\subsection{Economic Feasibility and Life Cycle Cost Analysis}

\noindent\textbf{Economic:} No dedicated hardware budget; use student devices and free/open map tiles (e.g., OpenStreetMap). Engineering cost is primarily time.

\begin{itemize}[leftmargin=1.5em]
    \item \textbf{Material/component cost:} \$0 incremental hardware (smartphones/laptops already owned by students). Optional accessories (power bank, phone holder) are non-essential and can be avoided.
    \item \textbf{Implementation cost:} The main cost driver is engineering effort. For a small team prototype, we assume on the order of 200--300 total person-hours across ENS~491/492 for coding, iteration, and reporting. This is an academic cost rather than a cash expense.
    \item \textbf{Operational cost:} For demos, static hosting (GitHub Pages) is \$0. If a server is needed later (not required for the current prototype), a minimal cloud instance could be \$5--\$15/month.
    \item \textbf{Maintenance and service cost:} Low for a static web prototype; periodic updates may be needed when browser APIs or libraries change (e.g., Leaflet/D3 upgrades). We mitigate this with pinned vendor assets and reproducible documentation.
    \item \textbf{End-of-life cost:} Minimal; the system can be archived as a static repository. Data retention policies (deleting sensitive traces, keeping aggregated results) are the primary ``lifecycle'' concern.
    \item \textbf{Cost per unit and performance-to-cost ratio:} For the functional unit (one analysis session), the marginal cash cost is effectively \$0, while the expected performance benefit is improved interpretability of location uncertainty and clearer error/deviation communication. Relative to conventional point-only displays, the value proposition is favorable because it improves decision quality without requiring additional infrastructure.
\end{itemize}

\section{RISK MANAGEMENT}

\begin{itemize}[leftmargin=1.5em]
    \item \textbf{Resource risk:} If GNSS logging devices are unavailable, use simulator data or open trajectory datasets \cite{Zheng2015} for development; prioritize one reliable phone for ground truth collection.
    \item \textbf{Expertise gap:} If advanced statistics or ML are needed, we will use coursework, advisor consultation, and documented libraries rather than ad-hoc models.
    \item \textbf{Schedule slip:} If implementation or ethics review delays user studies, deliver a strong prototype and heuristic evaluation first; formal study moves to ENS 492.
    \item \textbf{Plan B for visualization:} If a proposed encoding fails usability pretests, fall back to literature-supported baselines (circles, fades) \cite{mckenzie2016assessing} while iterating on one novel variant at a time.
    \item \textbf{Hosting/build:} If GitHub Pages policies change, distribute the static HTML or containerize locally; bibliography and figures remain in-repo.
\end{itemize}

\section{PROJECT SCHEDULE}

The course requests a Gantt chart. Figure~\ref{fig:ens491_gantt} shows the indicative timeline and task leaders; Table~\ref{tab:schedule} lists the same plan in tabular form. The chart can be updated as weeks shift. \textit{Delays in evaluation push final write-up; buffer weeks are reserved before advisor milestones.}

\begin{figure}[htbp]
    \centering
    \includegraphics[width=\linewidth]{images/ens491_gantt.png}
    \caption{Indicative Gantt chart for ENS~491 (GPS visualization track): weeks W1--W13, one leader per task.}
    \label{fig:ens491_gantt}
\end{figure}

\begin{table}[h]
    \centering
    \caption{Indicative schedule. Task leader = single responsible student.}
    \label{tab:schedule}
    \small
    \begin{tabular}{@{}lllp{3.2cm}@{}}
        \toprule
        \textbf{Task} & \textbf{Start} & \textbf{Duration} & \textbf{Leader} \\
        \midrule
        Literature extension & Week 1 & 3 weeks & Mustafa Bozyel \\
        Prototype: modes \& UI polish & Week 2 & 5 weeks & Ömer Mert Özel \\
        Design concepts (figures/docs) & Week 4 & 4 weeks & Senih Kırmaç \\
        Pilot evaluation / heuristics & Week 8 & 4 weeks & Mustafa Bozyel \\
        Draft report \& presentation & Week 11 & 3 weeks & Senih Kırmaç \\
        \bottomrule
    \end{tabular}
\end{table}

\noindent\textbf{Delay mitigation:} Parallel tracks (literature + coding) reduce critical-path length; weekly syncs surface blockers early.

\section{ETHICAL ISSUES}

\textbf{Location privacy:} GNSS trajectories can reveal home, work, and habits \cite{Zheng2015,Sobral2019}. We collect data only with explicit participant consent, minimize shared identifiers, and aggregate or blur public screenshots if needed.

\textbf{Misleading precision:} A well-designed uncertainty UI could still be misused to imply scientific rigor without validation; we will document limitations and avoid overstating accuracy of visual metaphors.

\textbf{AI-assisted design:} Exploratory mockups may use AI-generated ideas; these are labeled as non-final and must not infringe third-party rights.

\textbf{Dual use:} Movement visualization could theoretically support surveillance; our academic focus is transparency and user understanding, not covert tracking.

\noindent If future phases involve human subjects beyond informal team testing, we will follow university ethics procedures.

\section{REFERENCES}

\bibliographystyle{IEEEtran}
\bibliography{references}

\end{document}
